\chapter*{General Conclusion and Perspectives}
\addcontentsline{toc}{chapter}{General Conclusion and Perspectives}
\thispagestyle{fancy}
\justifying

The design and development of a digital decision support system for pre-hospital triage at the SAMU of Burkina Faso, based on the Emergency Severity Index (ESI) protocol, represents a fundamental achievement in the modernization of our emergency medical infrastructure. Throughout this project, we have transformed a manual and subjective process into a standardized, efficient, and secure digital platform that addresses the real-world challenges faced by SAMU operators and supervisors.

\section*{Synthesis of Contributions}
By adopting the Two-Track Unified Process (2TUP) methodology, we successfully navigated both the complex clinical logic of the 5-level ESI algorithm and the technical demands of a modern full-stack web application. The developed system provides a robust solution for patient prioritization, integrating real-time ESI level calculation with comprehensive audit logging and reporting features. Our functional testing and performance analysis have demonstrated that this digital approach significantly reduces ambiguity in patient categorization, ultimately optimizing the deployment of mobile medical units (SMUR). Furthermore, the integration of advanced security measures and intuitive dashboards ensures that the SAMU triage portal is production-ready and capable of handling sensitive medical information with the highest standards of integrity.

\section*{Key Perspectives for Further Development}
While the current version of the system provides a solid foundation for digitized triage, several avenues for future research and development have been identified. From a technical standpoint, the implementation of a mobile-first application with offline capabilities would further enhance the system's utility in remote or low-connectivity areas of Burkina Faso. Operationally, the integration of the triage portal with hospital emergency departments could streamline the handoff process, ensuring that receiving facilities are informed of incoming patients' priority levels in real time. Additionally, the potential for incorporating geographic information systems (GIS) for fleet tracking and optimization offers a promising direction for improving overall emergency response times.

Ultimately, the digital SAMU triage system is more than just a software tool; it is a clinical and operational asset that contributes to the core mission of saving lives by ensuring that the right care reaches the right patient at the right time.

\clearpage
